\scanpage{96}
\mspage{83}

que $\unsure{u|v}_{\unsure{\cdot|{}}}^*u_!:M(I_0)\longrightarrow M(I_0)$ \deleted{\unsure{transforme}} \inserted{applique} $C$ dans $C$, $\mathrm{TC}$ dans $\mathrm{TC}$. Or on a
\[
u_!(F_0)(i)=\varinjlim_{(\alpha,\deleted{\unsure{i}}j)\in i\backslash I_0}F_0(j)
\]
\margin{$\vert$ où \deleted{\unsure{$\alpha:u(i)\longrightarrow j=j$}} \inserted{$j\in\mathrm{Ob}\,I_0=\deleted{\unsure{\mathrm{Ob}}}$} $\mathrm{Ob}\,I$, et \deleted{$\alpha:i\longrightarrow$} $\alpha:i\longrightarrow u(j)=j$} et $i\backslash I_0$ est la catégorie discrète \inserted{au dessus de $I_0$}, dont la fibre en $j\in\mathrm{Ob}\,I_0=\mathrm{Ob}\,I$ est $\mathrm{Hom}_I(i,j)$ :
\[
i\backslash I_0=\coprod_{j\in\mathrm{Ob}\,I}\binom{\text{cat. discrète associée}}{\text{à }\mathrm{Hom}_I(i,j)}\,,
\]
donc
\diagram{\begin{tikzcd}[ampersand replacement=\&,column sep=small,row sep=small] u_!(F_0)(i)={} \& \displaystyle\coprod_{j\in\underset{\substack{\Vert\\\mathrm{Ob}\,I_0}}{\mathrm{Ob}\,I}} \underbrace{F_0(j)^{(\mathrm{Hom}_I(i,j))}}_{\substack{\text{somme directe}\\\text{de }\mathrm{Hom}_I(i,j)\\\text{copies de }F_0(j)}} ,\\ \text{\deleted{\unsure{donc}} on trouve} \& {}\\ \text{et}\quad u^*u_!(F_0)(\underset{\substack{\rotatebox[origin=c]{-90}{\(\in\)}\\\mathrm{Ob}\,I_0=\mathrm{Ob}\,I}}{i})=u_!(F_0)(i)={} \& {} \arrow[uu,bend right=23]
\end{tikzcd}\,,}

Comme $C\deleted{\unsure{(I_0)}}$ et $\mathrm{TC}\deleted{\unsure{(I_0)}}$ sont stables par sommes directes quelconques (l'étant par sommes directes finies, à cause des conditions de stabilité \unsure{finie|finies} sur $C$ et $\mathrm{TC}$, et puis par sommes infinies à cause de l'axiome \margin{\unsure{mérite} un lemme général \ldots}
% slanted
de stabilité transfinie \ldots), il s'ensuit \inserted{bien} que $u^*u_!$ \deleted{\unsure{transforme}} \inserted{applique} $C(I_0)$ et $\mathrm{TC}(I_0)$ en eux-mêmes,

Ainsi, la \deleted{\unsure{conclusion}} première assertion de 4.8. résulte de 4.6. De plus (4.\unsure{6|8}.8.) et (4.8.\unsure{4|9}) ont été vus dans 4.7., dans les conditions générales de 4.6 : et en supposant de plus le $4^\circ$ de (4.7.4)\unsure{)|{}}.
% archive sheet number: 95
